\emph{Setting.} The hypothesis is the one-loop formula of Gildener and Weinberg (Fact \claimref{synthesis.imported-coleman-weinberg-gildener-weinberg-scalon-mass}): when a scale-free tree potential has a flat direction $\varphi\hat{n}$ on which it vanishes, the one-loop potential along it is $B\varphi^{4}(\ln(\varphi^{2}/\mu^{2}) + \text{const})$ with $64\pi^{2}B v^{4} = \sum_{i}n_{i}M_{i}^{4}$ (vectors $n = 3$, real scalars $n = 1$, Dirac fermions $n = -4N_{c}$, masses at $\varphi = v$), and its minimum gives the scalon mass $m_{S}^{2} = 8Bv^{2}$, which is the displayed formula. It requires every field-dependent mass squared to be proportional to $\varphi^{2}$ and no explicit dimensionful coupling to the scalon; hypothesis (D) secures that for the series' content, and the paragraph \emph{Decoupling} below says why it is a restriction.

\emph{(i).} By \claimref{synthesis.rung-field-on-the-lift-locus}(2) the gauge mass form at a charge-preserving point is $\tfrac14 N P$ with $P$ of eigenvalues $g^{2}, g^{2}, g^{2} + g'^{2}, 0$, giving two vectors of mass $M_{W}$, one of mass $M_{Z}$ and the photon, all with $M^{2} \propto N$; with three polarisations each, $6M_{W}^{4} + 3M_{Z}^{4}$. By clause (3) of the same result, the six real directions on the locus are three gauge directions (eaten, massless in Landau gauge, no contribution), the neutral mode (the scalon itself, absent from the sum) and a charged plane of two real directions. On the flat physical ray of a scale-free potential (\claimref{synthesis.scale-free-rung-potentials-need-the-flat-physical-branch}, clause (3)) the charged mass squared is $2(2a - e)N$: along $\alpha = (0, \varepsilon z)$, which keeps $Q = \alpha^{\mathsf T}\beta = 0$ exactly, $N$ and $m$ both grow by $\tfrac12\varepsilon^{2}|z|^{2}$, and $V = aN^{2} + bm^{2} + eNm$ grows by $(2a - e)N\varepsilon^{2}|z|^{2} + \tfrac14(a + b + e)\varepsilon^{4}|z|^{4}$ at $m = -N$ with $b = e - a$; the kinetic coefficient is $|z|^{2}$. Under the conjugate reading every tree Dirac entry is $\lambda = |\beta| = \sqrt{2N}$ on the physical branch (clause (5)), so the fermion masses squared are proportional to $N$ as well; this holds in the per-sector reading and in its su(2)$_{R}$-covariant completion alike, since on the ray $\alpha = 0$ both $c$ and $c^{*}$ have size $|\beta|$. The top quark is a Dirac fermion with three colours, weight $-12$; the bottom and charm likewise, and the tau $-4$.

\emph{(ii) to (iv).} Exact rational arithmetic with the stored masses and brackets: $\pi \in [\tfrac{333}{106}, \tfrac{355}{113}]$; $v \in [246, 247]$ GeV, which contains $246.22$ GeV (the Fermi scale, \claimref{quantity.fermi-scale}) and the tree value $v^{2} = M_{W}^{2}\sin^{2}\theta_{W}\alpha^{-1}/\pi$ from the stored $M_{W}$, $\sin^{2}\theta_{W}(M_{Z})$ and $\alpha^{-1}(M_{Z})$, about $246.6$ GeV; $m_{t} \in [162, 173]$ GeV, which contains the stored pole mass $172.57$ GeV, the stored $m_{t}(M_{Z}) = 171.7$ GeV and the $\overline{\mathrm{MS}}$ value $m_{t}(m_{t})$ near $162.5$ GeV (from memory); $m_{h} \in [125, 125.4]$ GeV containing the stored $125.20$ GeV. Solving (i) for $m_{\pm}^{4}$ gives an expression increasing in $\pi$, $v$, $m_{h}$, $m_{t}$ and decreasing in $M_{W}$, $M_{Z}$, so its range over the box is attained at the corners, which the check evaluates exactly: $450^{4} < m_{\pm}^{4} < 456^{4}$. The threshold (iii) is the zero of the right side of (i) at $m_{h} = 0$: the check evaluates its sign at $249, 250$ GeV ($m_{t} = 162$) and $267, 268$ GeV ($m_{t} = 173$). The quartic: $m_{\pm}^{2} = 2(2a - e)N = (2a - e)v^{2}$, and the corners give $3.2 < 2a - e < 3.5$. Because the mass of the charged pair exceeds $m_{t} - m_{b}$ for every top mass in the range, $t \to bH^{+}$ is closed. The tree prediction $m_{\pm}^{2} = M_{W}^{2} + m_{h}^{2} - M_{Z}^{2}$ of $W = \kappa N^{2}$ (\claimref{synthesis.rung-field-d-term-action-and-tree-masses-on-the-lift-locus}(4) with \claimref{synthesis.rung-field-charged-scalar-sum-rule}) is not a flat-direction potential: its flat point has $m_{\pm}^{2} = M_{W}^{2} - M_{Z}^{2} < 0$ (\claimref{synthesis.scale-free-rung-potentials-need-the-flat-physical-branch}, clause (4)). The quartic $2a - e \approx 3.4$ is about $1.7$ in the Standard Model normalisation $m_{h}^{2} = 2\lambda v^{2}$: large; its running, and with it the scale of a Landau pole, is not computed here.

\emph{(v).} On the flat ray $m_{\pm}^{2} = (2a - e)\varphi^{2}$, so the scalon-pair coupling is $\partial m_{\pm}^{2}/\partial\varphi = 2m_{\pm}^{2}/v$ and the pair's term in the diphoton amplitude, $(g_{hSS}v/2m_{S}^{2})Q^{2}A_{0}$, has unit weight with $Q = 1$. The loop functions (cited from memory) are $A_{1} = -(2\tau^{2} + 3\tau + 3(2\tau - 1)f)/\tau^{2}$, $A_{1/2} = 2(\tau + (\tau - 1)f)/\tau^{2}$, $A_{0} = -(\tau - f)/\tau^{2}$, $f = \arcsin^{2}\sqrt{\tau}$ for $\tau \le 1$, and the top enters with $N_{c}Q_{t}^{2} = \tfrac43$. The check bounds $f$ by its series $\tfrac12\sum_{n}(4\tau)^{n}/(n^{2}\binom{2n}{n})$, whose terms are at most $\tau^{n}$, with a rational tail bound and outward rounding, verifies the bracket against $\arcsin^{2}\tfrac12 = \pi^{2}/36$, and obtains $-8.332 < A_{1}(\tau_{W}) < -8.331$, $1.835 < \tfrac43 A_{1/2}(\tau_{t}) < 1.836$, $0.336 < A_{0}(\tau_{\pm}) < 0.337$ and $0.898 < \mu_{\gamma\gamma} < 0.900$ for $m_{\pm} \in [450, 456]$ GeV. The measured diphoton signal strength, $1.10 \pm 0.06$ in the Particle Data Group average of the ATLAS and CMS measurements, lies about three standard deviations above; the ATLAS full Run 2 value alone, $1.04^{+0.10}_{-0.09}$ (arXiv:2207.00348), lies about 1.5 standard deviations above. Both values were checked against the published sources.

\emph{Decoupling.} Hypothesis (D) is not a property of the stored content. Under the conjugate reading the neutrino Dirac entries are read like the up-type ones (\claimref{synthesis.rung-field-on-the-lift-locus}(5)), so order-one Dirac couplings $y$ are allowed for the heaviest family, and the Pati-Salam route stores $v_{R} = M_{I}/g_{R} = 4.15 \times 10^{10}$ GeV (\claimref{standard-model.what-angle-asks-neutrino-sector}). Heavy Majorana neutrinos of mass $M_{R}$ then contribute to the rung field a one-loop mass term of order $y^{2}M_{R}^{2}/(4\pi^{2})$, an explicit dimensionful coupling far above $v^{2}$ that lifts the flat direction; requiring it below $v^{2}$ is the Vissani-type naturalness bound (F. Vissani, Phys. Rev. D 57 (1998) 7027, cited from memory, not checked), $M_{R}$ below about $10^{7}$ GeV for order-one $y$. The right-handed gauge bosons of the Pati-Salam stage couple to the rung bidoublet and give the same kind of term, of order $g_{R}^{2}M_{W_{R}}^{2}/(16\pi^{2})$, at one loop. Classical scale invariance removes no physical threshold, so (D) holds only if these couplings are absent or the heavy scales are near the weak scale; argued, not computed.

\emph{Remarks (computed, floating point, one loop, $\overline{\mathrm{MS}}$, Landau gauge; not part of the claim).} (a) Per-sector reading (not su(2)$_{R}$ covariant): for a potential even under $c \mapsto c^{*}$ both rays are flat; in the per-sector reading the fermions read only $|\beta|$, so on the mirror ray they are nearly massless (under the su(2)$_{R}$-covariant completion their entries there are of order $|\alpha|$ and the comparison below does not apply), and with the bosonic content equal the mirror minimum is deeper (ratio of minimum energies about $1.75$): the conjugation-odd quartic $e$ must be positive for the physical branch to be the vacuum. (b) With $W = 0$ as a boundary condition the physical-branch coefficient equals $(g^{2} + g'^{2})/8 \approx 0.069$, while its slope near the flat point is $m_{h}^{2}/v^{2} \approx 0.26$ per e-fold: the pure D-term boundary sits about $0.27$ e-folds above the transmutation scale, and the charged quartic $2a - e \approx 3.4$ there cannot run from the D-term's $\tfrac14 g^{2} \approx 0.1$ in that interval, so a pure D-term at a high scale is not a boundary for this route. The transmutation scale is therefore an input, in place of the level; under the Henneaux-Teitelboim promotion it is an input of $v/M_{P}$ type, like the level itself (\claimref{synthesis.constant-d-term-level-fixes-weak-scale-against-planck-under-promotion}).
